\chapter{Measurement Tables}
\label{app:measurement-tables}

Chapter~\ref{chap:evaluation} shows six sets of measurements as charts, because what each of them
reports is a comparison or a trend that is easier to see than to read. This appendix gives the same
measurements as tables, so that every value can be read exactly and so that the statistics a chart
leaves out, such as percentiles and cache-hit medians, remain available. Each chart reads the same
result files as its table, and a test in the repository checks that the numbers in the two agree;
Section~\ref{app:reproduce} names the files.

\section{Time-to-First-Frame}

Table~\ref{tab:ttff-results} gives the distribution of each stage of the baseline measurement against the S3
bucket; Figure~\ref{fig:ttff-cdn} shows its medians.

\begin{table}[H]
\centering
\caption{Time-to-First-Frame measurement results}
\label{tab:ttff-results}
\small
\begin{tabularx}{\textwidth}{X r r r r r}
\toprule
\textbf{STAGE} & \textbf{MEAN} & \textbf{MEDIAN} & \textbf{P95} & \textbf{MIN} & \textbf{MAX} \\
\midrule
Master playlist    & 95.55  & 82.49  & 355.02 & 76.61 & 355.02 \\
Rendition playlist & 88.62  & 82.40  & 207.32 & 77.62 & 207.32 \\
First segment      & 108.76 & 87.04  & 322.20 & 78.41 & 322.20 \\
\midrule
\textbf{Total TTFF} & \textbf{292.93} & \textbf{251.52} & \textbf{884.54} & \textbf{242.23} & \textbf{884.54} \\
\bottomrule
\end{tabularx}
\end{table}

\noindent\textit{Note: values in milliseconds, $n=20$ samples, measured against a video served
directly from the S3 bucket, before the CloudFront distribution described in
Section~\ref{sec:challenges} was deployed. These
figures therefore represent origin latency without any CDN, and they serve here as the baseline
against which Section~\ref{sec:cdn-measurement} compares the deployed system. The mean
(292.93~ms) exceeds the median (251.52~ms) because one outlier sample (884.54~ms) inflates both
the mean and the P95/max columns, which coincide with the single-sample maximum at this sample
size.}

Table~\ref{tab:ttff-cloudfront} gives the same statistics for the measurement through CloudFront after the
price class correction, which Figure~\ref{fig:ttff-cdn} compares with the baseline.

\begin{table}[H]
\centering
\caption{Time-to-First-Frame through CloudFront after the price class correction}
\label{tab:ttff-cloudfront}
\small
\begin{tabularx}{\textwidth}{X r r r r r}
\toprule
\textbf{STAGE} & \textbf{MEAN} & \textbf{MEDIAN} & \textbf{P95} & \textbf{MIN} & \textbf{MAX} \\
\midrule
Master playlist    & 28.65  & 8.37   & 263.43  & 6.54  & 263.43 \\
Rendition playlist & 29.29  & 7.93   & 348.78  & 5.89  & 348.78 \\
First segment      & 112.75 & 87.36  & 487.31  & 68.87 & 487.31 \\
\midrule
\textbf{Total TTFF} & \textbf{170.69} & \textbf{110.32} & \textbf{1099.51} & \textbf{83.17} & \textbf{1099.51} \\
\midrule
\textit{Warm cache only} ($n=18$) & \textit{108.49} & \textit{107.57} & \textit{159.25} & \textit{83.17} & \textit{159.25} \\
\bottomrule
\end{tabularx}
\end{table}

\noindent\textit{Note: values in milliseconds, $n=20$ samples with Signed Cookies attached. Two
samples were edge cache misses
and eighteen were hits; the final row reports only the samples that hit at all three stages,
which is the steady state a viewer of an already-watched video experiences. The P95 and maximum of
the full sample coincide with the colder of the two misses, as at this sample size the 95th
percentile is the single slowest observation.}

Table~\ref{tab:ttff-multiregion} adds, to the medians of Figure~\ref{fig:ttff-regions}, the edge location that
served each client and the median over the samples that hit the cache.

\begin{table}[H]
\centering
\caption{Time-to-First-Frame by client location, before and after \texttt{PriceClass\_All}}
\label{tab:ttff-multiregion}
\small
\begin{tabularx}{\textwidth}{X l r l r r}
\toprule
& \multicolumn{2}{c}{\textbf{PRICECLASS\_200}} & \multicolumn{3}{c}{\textbf{PRICECLASS\_ALL}} \\
\cmidrule(lr){2-3} \cmidrule(lr){4-6}
\textbf{CLIENT LOCATION} & \textbf{EDGE} & \textbf{MEDIAN} & \textbf{EDGE} & \textbf{MEDIAN} & \textbf{WARM} \\
\midrule
N. Virginia, United States & IAD61 & 55.23   & n/a   & n/a   & n/a \\
Dublin, Ireland            & DUB2  & 55.91   & DUB2  & 55.68 & 55.68 \\
Singapore                  & SIN3  & 56.89   & SIN3  & 53.20 & 53.20 \\
Tokyo, Japan               & NRT12 & 59.68   & n/a   & n/a   & n/a \\
Sydney, Australia          & n/a   & n/a     & SYD3  & 52.30 & 43.36 \\
S\~{a}o Paulo, Brazil      & CPT51 & \textbf{1196.10} & GRU1 & \textbf{36.79} & \textbf{36.49} \\
\midrule
Ho Chi Minh City (consumer) & SGN50 & 110.32 & n/a   & n/a   & n/a \\
\bottomrule
\end{tabularx}
\end{table}

\noindent\textit{Note: values in milliseconds, $n=20$ per location and per configuration. WARM
reports the median over samples that hit the
edge cache at all three stages. Entries marked n/a were not measured in that configuration: the
first run predated the Sydney probe, and the second re-measured only the two locations that could
change plus two controls. The Ho Chi Minh City row is the local measurement of
Table~\ref{tab:ttff-cloudfront}, taken from a consumer connection rather than from inside an AWS
region, and is included for comparison rather than as a like-for-like sample.}

\section{Playback Quality}

Table~\ref{tab:qoe-multisource} gives all five measures for each of the nine combinations of source and network
profile, of which Figure~\ref{fig:qoe-sources} charts two.

\begin{table}[H]
\centering
\caption{Playback quality across three sources of differing visual complexity}
\label{tab:qoe-multisource}
\footnotesize
\begin{tabularx}{\textwidth}{X r r r}
\toprule
\textbf{METRIC} & \textbf{LOW} & \textbf{MEDIUM} & \textbf{HIGH} \\
\midrule
\multicolumn{4}{l}{\textit{Unthrottled}} \\
Rebuffering ratio, mean     & 0.00\%  & 0.00\%  & 0.00\%  \\
Rebuffering ratio, max      & 0.00\%  & 0.00\%  & 0.00\%  \\
Runs exhibiting a stall     & 0/5     & 0/5     & 0/5     \\
Startup delay, median (s)   & 1.46    & 1.34    & 1.21    \\
Average bitrate (kbit/s)    & 1644    & 1830    & 1664    \\
\midrule
\multicolumn{4}{l}{\textit{Fast 3G}} \\
Rebuffering ratio, mean     & 0.22\%  & 2.30\%  & 0.11\%  \\
Rebuffering ratio, max      & 0.60\%  & 5.81\%  & 0.55\%  \\
Runs exhibiting a stall     & 2/5     & 2/5     & 1/5     \\
Startup delay, median (s)   & 1.22    & 1.15    & 1.12    \\
Average bitrate (kbit/s)    & 1477    & 1316    & 1246    \\
\midrule
\multicolumn{4}{l}{\textit{Slow 3G}} \\
Rebuffering ratio, mean     & 4.91\%  & 5.80\%  & 4.84\%  \\
Rebuffering ratio, max      & 13.14\% & 13.73\% & 13.37\% \\
Runs exhibiting a stall     & 3/5     & 3/5     & 3/5     \\
Startup delay, median (s)   & 3.76    & 1.87    & 3.84    \\
Average bitrate (kbit/s)    & 1070    & 1224    & 845     \\
\bottomrule
\end{tabularx}
\end{table}

\noindent\textit{Note: five recorded runs per cell after two discarded warm-up runs, 120~s of
playback each, at a device pixel ratio of one. Source durations are 4:23, 4:23, and 4:40, all
longer than the measurement window.}

\section{Transcoding Latency and Capacity}

Table~\ref{tab:chunk-config} lists the ten runs of the 25-minute video with the columns that
Figure~\ref{fig:chunk-config} does not show: capacity, slots and the length of the chunk phase.

\begin{table}[H]
\centering
\caption{A 25-minute video on the chunked path, by chunk configuration and capacity}
\label{tab:chunk-config}
\small
% GENERATED by scripts/results/build-report-tables.py from docs/results. Do not edit by hand.
\begin{tabularx}{\textwidth}{l l >{\raggedleft\arraybackslash}X >{\raggedleft\arraybackslash}X >{\raggedleft\arraybackslash}X >{\raggedleft\arraybackslash}X >{\raggedleft\arraybackslash}X}
\toprule
\textbf{N $\times$ VCPU} & \textbf{FARGATE} & \textbf{VCPU} & \textbf{SLOTS} & \textbf{PHASE (MIN)} & \textbf{READY (MIN)} & \textbf{VCPU-S} \\
\midrule
5 $\times$ 1 & Spot & 8 & 8 & n/a & 30.1 & n/a \\
5 $\times$ 1$^{\dagger}$ & Spot & 8 & 8 & 32.6 & 36.3 & 7\,117 \\
23 $\times$ 1 & Spot & 8 & 8 & 19.2 & 23.6 & 7\,270 \\
12 $\times$ 2 & Spot & 8 & 4 & 16.3 & 20.1 & 6\,816 \\
6 $\times$ 4 & Spot & 8 & 2 & 16.8 & 20.4 & 6\,792 \\
6 $\times$ 4 & Spot & 8 & 2 & 14.4 & 18.0 & 6\,363 \\
6 $\times$ 4 & On-Demand & 8 & 2 & 21.7 & 26.0 & 10\,182 \\
6 $\times$ 4 & On-Demand & 8 & 2 & 18.2 & 22.5 & 7\,852 \\
12 $\times$ 4 & Spot + OD & 16 & 4 & 11.5 & 15.3 & 8\,827 \\
12 $\times$ 4 & Spot + OD & 16 & 4 & 10.2 & 14.7 & 8\,159 \\
\bottomrule
\end{tabularx}

\end{table}

\noindent\textit{Note: the first four rows were timed from the creation of the first job and the
later rows from the end of the upload, which differ by about half a minute. The vCPU-seconds of the
first five rows count chunk jobs only; the later rows count every job of the video. Rows are in
chronological order of measurement. The first column gives the number of chunks and the vCPU of each; $^{\dagger}$ marks the run on the production environment. PHASE is the time from the first chunk start to the last chunk end, SLOTS the
number of chunks encoded at once, VCPU-S the vCPU-seconds, and OD stands for On-Demand.}

Table~\ref{tab:four-hour-breakdown} gives the minutes of each step of the four-hour video, including the steps too
short to be labelled in Figure~\ref{fig:four-hour}, and the compute used.

\begin{table}[H]
\centering
\caption{A four-hour video at two capacities}
\label{tab:four-hour-breakdown}
\small
% GENERATED by scripts/results/build-report-tables.py from docs/results. Do not edit by hand.
\begin{tabularx}{\textwidth}{X r}
\toprule
\textbf{STEP} & \textbf{MINUTES} \\
\midrule
Planner (Lambda submit to child jobs) & 0.4 \\
Audio 64 kbit/s and 128 kbit/s (in parallel) & 3.6 / 10.2 \\
48 chunk jobs, 8 at a time, 1 vCPU each & 170.4 \\
Finalizer (playlists, thumbnail, moderation, READY) & 1.9 \\
\textbf{Upload complete to READY} & \textbf{185.0} \\
\bottomrule
\end{tabularx}

\end{table}

Table~\ref{tab:concurrent-four} gives, for each of the four videos, the start and end of the chunk phase that
Figure~\ref{fig:concurrent-four} draws.

\begin{table}[H]
\centering
\caption{Four 25-minute videos submitted together, by capacity}
\label{tab:concurrent-four}
\small
% GENERATED by scripts/results/build-report-tables.py from docs/results. Do not edit by hand.
\begin{tabularx}{\textwidth}{>{\raggedleft\arraybackslash}X >{\raggedleft\arraybackslash}X >{\raggedleft\arraybackslash}X >{\raggedleft\arraybackslash}X >{\raggedleft\arraybackslash}X}
\toprule
\textbf{VCPU} & \textbf{ORDER} & \textbf{START (MIN)} & \textbf{END (MIN)} & \textbf{READY (MIN)} \\
\midrule
8 & 1 & 2.8 & 19.7 & 21.0 \\
8 & 2 & 18.4 & 37.1 & 38.7 \\
8 & 3 & 35.5 & 58.5 & 62.9 \\
8 & 4 & 52.2 & 71.9 & 73.3 \\
16 & 1 & 3.0 & 16.0 & 17.7 \\
16 & 2 & 2.7 & 24.6 & 25.7 \\
16 & 3 & 22.2 & 35.0 & 36.1 \\
16 & 4 & 32.0 & 43.3 & 44.3 \\
\bottomrule
\end{tabularx}

\end{table}

\noindent\textit{Note: in the 8~vCPU rows the four videos were uploaded together over one
connection and timed from the first job creation. In the 16~vCPU rows the arrival was produced by
server-side copies of one uploaded file, so all four events reached the pipeline within a second. START and END are the first chunk start and the last chunk end, and ORDER is the order of completion.}
